%%%% RESUMO
%%
%% Apresentação concisa dos pontos relevantes de um texto, fornecendo uma visão rápida e clara do conteúdo e das conclusões do
%% trabalho.

\begin{resumoutfpr}
Este Trabalho de Conclusão de Curso apresentará o processo de refatoração, modernização e expansão do sistema web do E-Museu, uma plataforma digital destinada à catalogação e divulgação de peças da história da informática. A justificativa para a realização deste trabalho se fundamentará nas limitações técnicas identificadas na versão atual da aplicação, que apresenta dificuldades de manutenção, ausência de padronização no código, insuficiência de testes automatizados e carência de funcionalidades relevantes para a continuidade do projeto. Diante desse cenário, o trabalho terá como objetivo aprimorar a estrutura interna do sistema, aumentar sua estabilidade e sustentabilidade a longo prazo e, simultaneamente, implementar novas funcionalidades que ampliarão a utilidade prática da plataforma. Como metodologia, será adotada uma abordagem baseada em boas práticas de engenharia de software, incluindo a aplicação de padrões de projeto, a reorganização de módulos, a padronização de versionamento e a configuração de ambientes de desenvolvimento, teste e produção. Além disso, serão utilizados princípios de refatoração para melhorar legibilidade, modularidade e desempenho do código, enquanto testes automatizados serão incorporados para garantir a confiabilidade das funcionalidades e prevenir problemas futuros. Entre os principais resultados esperados estarão a reestruturação completa do código-fonte, a melhoria significativa na organização e rastreabilidade do projeto, a implementação de mecanismos de segurança, o aumento da capacidade de upload de mídias, a internacionalização do sistema e a geração de etiquetas para catalogação física das peças. Espera-se que a conclusão evidencie que o conjunto de melhorias aplicadas tornará o E-Museu mais robusto e coerente com práticas modernas de desenvolvimento, estabelecendo uma base sólida para evoluções futuras e contribuindo para a preservação da memória tecnológica.
\end{resumoutfpr}
